%  LaTeX support: latex@mdpi.com 
%  For support, please attach all files needed for compiling as well as the log file, and specify your operating system, LaTeX version, and LaTeX editor.

%=================================================================
\documentclass[geomatics,article,submit,pdftex,moreauthors]{Definitions/mdpi} 



% MDPI internal commands
\firstpage{1} 
\makeatletter 
\setcounter{page}{\@firstpage} 
\makeatother
\pubvolume{1}
\issuenum{1}
\articlenumber{0}
\pubyear{2023}
\copyrightyear{2023}
%\externaleditor{Academic Editor: Firstname Lastname}
\datereceived{16 July 2023} 
\daterevised{} 
\dateaccepted{} 
\datepublished{} 
\hreflink{https://doi.org/} % If needed use \linebreak
%\doinum{}

%=================================================================
% Add packages and commands here. The following packages are loaded in our class file: fontenc, inputenc, calc, indentfirst, fancyhdr, graphicx, epstopdf, lastpage, ifthen, lineno, float, amsmath, setspace, enumitem, mathpazo, booktabs, titlesec, etoolbox, tabto, xcolor, soul, multirow, microtype, tikz, totcount, changepage, attrib, upgreek, cleveref, amsthm, hyphenat, natbib, hyperref, footmisc, url, geometry, newfloat, caption

\graphicspath{{figures/}} % path to folder with figures
\usepackage{subcaption}
\usepackage{listings}
\usepackage{xcolor}
\usepackage{gensymb} % for \textdegree
\usepackage{tabularx}

%=================================================================
% Full title of the paper (Capitalized)
\Title{Seafloor and Ocean Crust Structure of the Kerguelen Plateau from Marine Geophysical and Satellite Altimetry Datasets}
% Interplay Between Seafloor Spreading, Marine Gravity Roughness and Ocean Crust Structure of the Kerguelen Plateau Analysed from Satellite Altimetry

% MDPI internal command: Title for citation in the left column
\TitleCitation{Seafloor and Ocean Crust Structure of the Kerguelen Plateau from Marine Geophysical and Satellite Altimetry Datasets}

% Author Orchid ID: enter ID or remove command
\newcommand{\orcidauthorA}{0000-0002-5759-1089} % Polina Add \orcidA{} behind the author's name

% Authors, for the paper (add full first names)
\Author{Polina Lemenkova $^{1}$*\orcidA{}}

%\longauthorlist{yes}

% MDPI internal command: Authors, for metadata in PDF
\AuthorNames{Polina Lemenkova}

% MDPI internal command: Authors, for citation in the left column
\AuthorCitation{Lemenkova, P.}
% If this is a Chicago style journal: Lastname, Firstname, Firstname Lastname, and Firstname Lastname.

% Affiliations / Addresses (Add [1] after \address if there is only one affiliation.)
\address{%
$^{1}$ \quad Laboratory of Image Synthesis and Analysis (LISA), École polytechnique de Bruxelles (Brussels Faculty of Engineering), Université Libre de Bruxelles (ULB). Building L, Campus du Solbosch, ULB -- LISA CP165/57, Avenue Franklin D. Roosevelt 50, Brussels 1050, Belgium.
}

% Contact information of the corresponding author
\corres{Correspondence: polina.lemenkova@ulb.be; Tel.: +32 471 86 04 59 (P.L.)}

% Current address and/or shared authorship
%\firstnote{Current address: Affiliation 3.} 

% Abstract (Do not insert blank lines, i.e. \\) 
\abstract{The volcanic Kerguelen Islands are formed on one of the world's largest submarine plateaus. Located in the remote segment of the southern Indian Ocean close to Antarctica, the Kerguelen Plateau is notable for complex tectonic origin and geologic formation related to the Cretaceous history of the continents. This is reflected in varying age of the oceanic crust adjacent to the plateau and highly heterogeneous bathymetry of the Keguelen Plateau where seafloor structure of the southern segment  differs significantly from its northern part. Remote sensing data derived from marine gravity and satellite radar altimetry surveys serve as an important source of information for mapping of such complex seafloor features. Automatic matching of topographic and geophysical grids with high accuracy is essential for complex geologic-tectonic investigations. We propose incorporating the geospatial information from NOAA, EMAG2 and WDMAM, ETOPO1, EGM96 and EGM2008 to refine the extent and distribution of features extracted from geospatial datasets. The cartographic joint analysis of topography, magnetic anomalies, tectonic and gravity data is based on integrated mapping using scripting toolset Generic Mapping Tools (GMT). The refined mapping of the submerged features (Broken Ridge orientation, seamounts around the Crozet Islands, fabric of seafloor and orientation and frequency of magnetic anomalies) is used to analyse the correspondence of the objects with regard to free-air gravity and magnetic anomalies, geodynamics and seabed structure in south-west Indian Ocean. Experimental results show that integrating the proposed datasets using the state-of-the-art cartographic scripting language improves identification and visualization of the seabed objects. The presented research includes 10 new maps which contribute to  increasing the knowledge of subsurface seafloor structure around the Kerguelen Plateau and composition of the Earth's crust in the French Southern and Antarctic Lands.
}

% Keywords
\keyword{Antarctic; Southern Ocean; bathymetry; French Southern and Antarctic Lands; cartography; satellite altimetry; marine geophysics; keyword 8; keyword 9; keyword 10} 

% The fields PACS, MSC, and JEL may be left empty or commented out if not applicable
%\PACS{J0101}
%\MSC{}
%\JEL{}

\begin{document}

%----------- section ----------->
\section{Introduction}

%----------- subsection ----------->
\subsection{Background}
Analysis of seafloor structure and ocean crust concerns mapping, detecting, and recognising the bathymetric structures and variations in geophysical fields analysed from the remote sensing data. Among these problems, integrated geophysics interpretation of satellite altimetry,  marine gravimetry and magnetic anomalies is one of the most active research areas and has attracted much research interest over the past decades. Previous studies \cite{ESHAGH2021313} provided a broad overview of numerous approaches for analysing satellite gravimetry data to model the Earth beneath the sea. Marine geophysical and satellite altimetry data provide information on the crustal density and processes in the upper mantle \cite{KAPAWAR2021106692,PENG2022110823,ZHANG202368}, sediment thickness and basement depths \cite{Lemenkova202005}. Moreover, the analysis of the geophysical data enables to reveal the key characteristics regarding the lithosphere thickness such as Moho discontinuity, elastic thickness \cite{TESAURO201378}, viscosity and flexural rigidity. Such data are essential for geological investigations and our understanding of the Earth's structure.

In terms of the data-driven analyses, previous works investigating the oceanic seabed can be roughly classified into studies focused on bathymetric mapping and geophysical analysis. Mapping seafloor usually characterises the bathymetric patterns with increased cartographic details of data visualization and improved methods \cite{Battista,Lemenkova202105,Lemenkova202012}. Marine geophysical methods generally investigate the geology of continental margins with deep-sea drilling \cite{Schlich1992TheGA,Veevers1991a}, analysis of features extracted from seismic survey and remote sensing data \cite{NEGI2022229330}, investigating the structure of the deep ocean basins \cite{BENARD2010633,RAMANA2001241,DAVIS201616} and modelling mid-ocean ridge system in a global or regional context \cite{FREY200073,Sreejith,Lemenkova202112}. However, methods of bathymetric survey are of comparatively high complexity and high cost, as they require detailed hydrographic equipment such as mutlibeam echo sounder systems onboard the research cruise \cite{Falvey}. Although they supply novel information on origin, evolution, structure, stratigraphy, and tectonic features of the oceanic crust, they also do not provide matching of the multi-source geophysical data for modelling seafloor in an explicit way. 

%----------- subsection ----------->
\subsection{Related work}

It is known that seafloor bathymetry and geophysical setting of the oceanic crust are two important cues for modelling lithosphere \cite{WATSON201697}. However, the integrated capture and processing of multi-source data, -- such as bathymetry, marine geologic data (development on the Earth's crust including age, spreading rates and symmetry), marine geophysical data (sediment thickness, gravity and magnetic anomalies), -- has always been a challenging task, and few attempts have been made to explore it. An integrated use of the satellite-derived data from Cryosat-2 and Saral/AltiKa satellite altimetry and the ship-measured gravity data were applied for for estimating marine gravity anomalies \cite{NGUYEN2020505}. Mapping lithospheric structures such as tectonic blocks is performed based on the satellite gravity data including World Gravity Model (WGM) and Gravity Recovery and Climate Experiment (GRACE) \cite{EMISHAW2023104775}. The Earth Gravity Model (EGM) and astrogeodetic vertical deflections were employed for modelling gravity  and geopotential difference \cite{WANG2023}. A joint processing of terrestrial and satellite geodetic data was applied for solving theoretical issues of geodesy such as altimetry–gravimetry boundary-value problem \cite{Lehmann}.

The use of remote sensing data, satellite altimetry and gravimetry, has proven to be effective in geophysical and environmental studies. Numerous examples exist in recent literature with case studies on estimating geoid values, evaluating gravity anomalies, tectonic and crustal structures, glacial and hydrological research and habitat mapping. Thus, satellite altimetry data can be used to jointly represent patterns of oceanic currents and draw conclusions on connectivity between habitats \cite{MORI201645}. A study by \cite{COTTE2022103650} is in part similar in using hydrographic and acoustic data for analysis of vertical layers of the ocean to identify community distribution and carbon content. Tide modelling with regard to bathymetric gradients was computed based on satellite altimetry data \cite{Maraldi}. The retreat of glaciers using mass balance measurements was estimated for analysis of spatial distribution of surface mass balance \cite{Verfaillie}. Mathieu et al. \cite{MATHIEU2011206} combine SRTM DEM, remote sensing data and geological field campaign to clarify the structures of Kerguelen Islands. 

Such examples illustrate that a combination of bathymetric and geophysical data for detailed analysis of the Earth's crust structure and topography outperforms their use separately, and suggests that seafloor mapping should be based on the analysis of the heterogeneous features of the seabed derived from various datasets, as pointed earlier \cite{COWLEY2004267,Lemenkova202125,1993A234,ABUAMARAH2019868}. Since such methods can be regarded as feature-level combinations of seafloor fabric showing local appearance of prominent features such as fracture zones, ridges, seamounts, deep-sea trenches, canyons and rifts, the use of scripting programming approach to fuse different integrated geophysical datasets enables to match data from various sources, resolution and origin for a detailed analysis of the seafloor structure and geodynamic setting \cite{Veevers1991b,Lemenkova202104,Hill}. More applications of the use of remote sensing data on Kerguelen focus on marine biology which is explained by a unique wildlife with vulnerable species, well-preserved due to the remote isolated location of the archipelago \cite{fishes8040174}, and seasonal effects of the circumpolar currents on the fish communities \cite{Damerell,MONGIN2008880,Park}. A combination of such factors makes the ecosystems of Kerguelen an environmental heritage \cite{Sombetzkib,Sombetzkia,Naim}. 

%----------- subsection ----------->
\subsection{Objection and motivation}

In our approach, we combine the bathymetric and marine geophysical data on Kerguelen Plateau in regional and local scales with varied level of details to describe the structures of the south-west Indian Ocean. Specifically, the presented study aims at investigating the geophysical and topographic setting of the Kerguelen Plateau using the integrated datasets which cannot be achieved with separate use of data. Ten new maps are presented and described to visualise spatial variations, reveal correlation and capture the matches between geophysical and topographic setting of the selected segment of the south-west Indian Ocean that belongs to the French Southern and Antarctic Lands. Given the remote location of the Kerguelen Archipelago, one of the most isolated places on Earth, and high cost of geological and geophysical sampling, the integrated use of the satellite derived data presented in this study enables to get better insights into the structure of the Kerguelen Plateau and south-west Indian Ocean.

%----------- section ----------->
\section{Study area}

The Kerguelen Plateau presents an extended topographic elevation situation in the south-western sector of the Indian Ocean, Figure \ref{fig01}. The structure of its major geographic boundaries are as follows. To the northwest, it is limited by the Crozet Archipelago located WNW 1250 km \cite{BRETON2013126}, to the north -- by the African-Antarctic Basins, to the northeast -- by the Australian-Antarctic Basin, and to the southwest -- by the Enderby Basin, Figure \ref{fig01}. 

\begin{figure}[H]
\makebox[1.0\linewidth][r]{%
\centering
	\hspace{100pt}\includegraphics[width=14.0 cm]{Fig_01.jpg}
}
\caption{Map of the Kerguelen Plateau region. Mapping: GMT. Map source: author. \label{fig01}}
\end{figure}

The origin of the Kerguelen Plateau is related to the hotspot associated with the the kinematics of Gondwana breakup and spreading between India and Antarctica \cite{Guglielmi}. The early surface expression of the Kerguelen Plume at the Southern Kerguelen Plateau is estimated in Mesozoic around 120 Ma \cite{Gaina}. Currently, it presents one of the world's largest the volcanic plateaus (the total length of 2300 km) with occasional evidence of volcanism remaining on the Heard and McDonald Islands. \cite{Li}. Morphologically, Kerguelen Plateau is oriented in a NNW direction toward the Antarctic continental margin where it is constrained by the Princess Elizabeth Land. Kerguelen Plateau is asymmetric in its structure, with notable variations in two parts. Northern part of the Kerguelen Archipelago (Heard, McDonald and Kerguelen Islands) is more shallow (<1000 m) compared to the southern segment (depths of 1500-2000 m) \cite{Ramsay}. 

%----------- section ----------->
\section{Materials and Methods}

%----------- subsection ----------->
\subsection{Data}

The materials used in this study as raw cartographic grids include the following datasets. The bathymetric mapping was based on the ETOPO Global Relief Model \cite{NOAA}; the gravity grids were derived from EGM-2008 \cite{Pavlis}. The data on sediment thickness are derived from the GlobSed dataset showing the distribution of the sediment thickness in the World's Oceans \cite{Straume}. The magnetic data are derived from the two available information sources compiled originally from satellite, airborne, and marine magnetic measurements: the Earth Magnetic Anomaly Grid (EMAG2) \cite{Maus}, which has a resolution of 2 arc minutes, and the World Digital Magnetic Anomaly Map (WDMAM) which has 3-minute resolution \cite{Lesur}. The marine geological data regarding age, spreading rates, and asymmetry of the ocean crust were derived from the dataset originally developed by Müller et al. \cite{Mueller}. 

%----------- subsection ----------->
\subsection{Methods}

 In this study, the cartographic methodology is based on using the Generic Mapping Tools (GMT) programming suite version 6-1-1 \cite{Wessel} using developed scripting workflow \cite{Lemenkova202212,Lemenkova202203}. The most prominent feature of GMT is using scripting approach that principally differs it from the conventional software due to the embedded programming language. Major advantages of GMT include modular toolbox system with over 140 existing modules which enables to generate highly customised mapping workflows, high print-quality graphics, geodetic-precision projections, reliable data processing algorithms applicable to large datasets and raster grids, flexibility to operating systems and open source licence \cite{WesselLuis}. 

While traditional GIS-based methods usually employ graphical user interface (GUI) for mapping and data visualization with existing standards menu of the software, script-based methods of GMT capture the dataset by running a script written using embedded language which operates similar to the programming approaches. Script consists of a consequence of several commands with pre-defined parameters that control the appearance of the cartographic elements on a map. The most well-known cartographic tools that use scripts for data processing are GMT \cite{Lemenkova202217} and GRASS GIS \cite{NETELER2012124}. While the latter is mostly used for image processing, theoretical cartography and ecological studies with numerous existing examples  \cite{ROCCHINI2017166,FRIGERI20111265,Lemenkova202313,SOROKINE2007685}, GMT is applied in geophysical and geological studies \cite{Lee,Lemenkova202206,DeSanto}.

%----------- section ----------->
\pagebreak
\section{Results and Discussion}

%----------- subsection ----------->
\subsection{Ocean floor formation}

Age, spreading rates and spreading symmetry of the ocean crust indicate the evolution steps in ocean floor formation in South-West Indian Ocean and the Kerguelen Plateau. The age of the ocean crust was determined by interpolation of adjacent seafloor isochrons oriented towards the direction of seafloor spreading, Figure \ref{fig02}.

\begin{figure}[H]
\makebox[1.0\linewidth][r]{%
\centering
	\hspace{100pt}\includegraphics[width=14.0 cm]{Fig_02.jpg}
}
\caption{Age of the ocean crust in SW Indian Ocean and the Kerguelen Plateau region and east Antarctic. Mapping software: GMT v. 6-1-1. Map source: author. \label{fig02}}
\end{figure}


\begin{figure}[H]
\makebox[1.0\linewidth][r]{%
\centering
	\hspace{100pt}\includegraphics[width=14.0 cm]{Fig_03.jpg}
}
\caption{Age, spreading rates and spreading symmetry of the ocean crust over the Kerguelen Plateau region, east Antarctic and south-west Indian Ocean: asymmetries in crustal accretion (\%) on conjugate ridge flanks. Mapping software: GMT v. 6-1-1. Map source: author. \label{fig03}}
\end{figure}

Spreading half rates of the ocean crust are visualised based on the 2 arc min netCDF grid v.3.6 (NOAA). 

\begin{figure}[H]
\makebox[1.0\linewidth][r]{%
\centering
	\hspace{100pt}\includegraphics[width=14.0 cm]{Fig_04.jpg}
}
\caption{Spreading half rates of the ocean crust over the Kerguelen Plateau region, east Antarctic and south-west Indian Ocean with added GSFML data. Mapping software: GMT v. 6-1-1. Map source: author. \label{fig04}}
\end{figure}

\begin{figure}[H]
\makebox[1.0\linewidth][r]{%
\centering
	\hspace{100pt}\includegraphics[width=14.0 cm]{Fig_05.jpg}
}
\caption{Age uncertainty in lithosphere crust (m.y) over the Kerguelen Plateau region, east Antarctic and south-west Indian Ocean. Cartography: GMT, version 6-1-1. Map source: author. \label{fig05}}
\end{figure}

%----------- subsection ----------->
\subsection{Sediment thickness}

\begin{figure}[H]
\makebox[1.0\linewidth][r]{%
\centering
	\hspace{100pt}\includegraphics[width=14.0 cm]{Fig_06.jpg}
}
\caption{Sediment thickness (m) in East Antarctica, South-West Indian Ocean and the Kerguelen Plateau region. Cartography: GMT. Lambert azimuthal equal-area projection. Map source: author. \label{fig06}}
\end{figure}

%----------- subsection ----------->
\subsection{Free-air gravity anomaly and vertical gravity gradients}

\begin{figure}[H]
\makebox[1.0\linewidth][r]{%
\centering
	\hspace{100pt}\includegraphics[width=14.0 cm]{Fig_07.jpg}
}
\caption{Marine gravity field over the Kerguelen Plateau region, south-west Indian Ocean. Mapping software: GMT v. 6-1-1. Map source: author. \label{fig07}}
\end{figure}

\begin{figure}[H]
\makebox[1.0\linewidth][r]{%
\centering
	\hspace{100pt}\includegraphics[width=14.0 cm]{Fig_08.jpg}
}
\caption{Vertical gravity gradient anomalies over the Kerguelen Plateau region showing variations in gravity with the change in topographic elevations. Mapping software: GMT v. 6-1-1. Map source: author. \label{fig08}}
\end{figure}

%----------- subsection ----------->
\subsection{Magnetic anomalies over Kerguelen Plateau}

The anomaly of the magnetic intensity at an altitude of 5 km above mean sea level over Kerguelen is shown in Figure \ref{fig09}. The analysis of the magnetic anomaly patterns in the south-west Indian Ocean supports the hypothesis of the spreading seafloor with variations in the oceanic crustal blocks movements, which is visible on different magnetic pattern over the mid-ocean ridge, as reported earlier \cite{LePichon}.

\begin{figure}[H]
\makebox[1.0\linewidth][r]{%
	\begin{subfigure}[b]{.65\textwidth}
		\centering
			\includegraphics[width=9.0cm]{Fig_09a.jpg}
		\caption{}
	\end{subfigure}	
	%\vspace{1mm}
	\begin{subfigure}[b]{.60\textwidth}
		\centering
			\includegraphics[width=9.5cm]{Fig_09b.jpg}
		\caption{}
	\end{subfigure}%
}
\caption{(\textbf{a}): Marine and Earth airborne magnetic anomaly grid (3-min resolution) based on WDMAM, south-west Indian Ocean; (\textbf{b}): Marine and Earth airborne magnetic anomaly grid based on EMAG-2 over the Kerguelen Plateau region. Black areas signify "no data" in the original grid. Cartography: GMT, version 6-1-1. Map source: author.}\label{fig09}
\end{figure}

%----------- subsection ----------->
\subsection{Geoid models}
The variations in the geopotential geoid models over Kerguelen is shown in Figure \ref{fig10}. The high anomalies in geoid level above the Kerguelen Plateau are explained by the high volcanism associated to a ridge formation on hot lithosphere, which is reflected in the extraordinary thickness below Kerguelen with geoid values of above 40 m (red areas in Figure \ref{fig10}) which exceeds normal thickness of the oceanic crust. This proves previous investigations on geoid in the Kerguelen Archipelago that reported anomalous thickness in this area \cite{RECQ1987301}. Such isostatic compensation associated with the anomalously high geoid values is related to the topographically rough elevated surfaces with regions of active tectonic uplift signifying upper mantle processes \cite{Coblentz}. Comparison of the geoid map with topography (sf. Figures \ref{fig01} and \ref{fig10}) implies correlation between the continued geodynamic processes in south-west Indian Ocean and the topography structure of the Kerguelen Archipelago. Moreover, this proves a high positive correlation between geoid height and deep structure of the seafloor topography, which is noted earlier for the regions of large plateaus and swells \cite{SandwellMacKenzie}. 

\begin{figure}[H]
\makebox[1.0\linewidth][r]{%
	\begin{subfigure}[b]{.75\textwidth}
		\centering
			\includegraphics[width=9.6cm]{Fig_10a.jpg}
		\caption{}
	\end{subfigure}	\\
	
	\begin{subfigure}[b]{.60\textwidth}
		\centering
			\includegraphics[width=9.5cm]{Fig_10b.jpg}
		\caption{}
	\end{subfigure}%
}
\caption{(\textbf{a}): Geoid model based on EGM-96 over the Kerguelen Plateau region, east Antarctic and south-west Indian Ocean; (\textbf{b}): Geoid model based on EGM-2008 over the Kerguelen Plateau region, east Antarctic and south-west Indian Ocean. Cartography: GMT, version 6-1-1. Map source: author.}\label{fig10}
\end{figure}

%----------- section ----------->
\section{Conclusions}

%----------- section ----------->
\authorcontributions{Supervision, conceptualization, methodology, software, resources, funding acquisition, and project administration, writing---original draft preparation, methodology, software, data curation, visualization, formal analysis, validation, writing---review and editing, and investigation, P.L.}

\funding{The publication was funded by the Editorial Office of \emph{Land}, Multidisciplinary Digital Publishing Institute (MDPI), by providing 100\% discount for the APC of this manuscript.}

\institutionalreview{Not applicable.}

\informedconsent{Not applicable.}

\dataavailability{Not applicable} 

\acknowledgments{The authors thank the reviewers for reading and review of this manuscript.}

\conflictsofinterest{The authors declare no conflict of interest.} 

\abbreviations{Abbreviations}{
The following abbreviations are used in this manuscript:\\

\noindent 
\begin{tabular}{@{}ll}
DEM & Digital Elevation Model \\
EGM2008 & Earth Gravitational Model 2008\\
EMAG-2 & Earth Magnetic Anomaly Grid (EMAG) 2 \\
GRACE & Gravity Recovery and Climate Experiment \\
GMT & Generic Mapping Tools \\
NOAA & National Oceanic and Atmospheric Administration \\
SRTM & Shuttle Radar Topography Mission \\
WGM & World Gravity Model \\
WDMAM & World Digital Magnetic Anomaly Map\\
\end{tabular}
}

%%%%%%%%%%%%%%%%%%%%%%%%%%%%%%%%%%%%%%%%%%
\begin{adjustwidth}{-\extralength}{0cm}
%\printendnotes[custom] % Un-comment to print a list of endnotes

\reftitle{References}
%\nocite{*}

%=====================================
% References, variant A: external bibliography
%=====================================
\bibliography{Kerguelen.bib}

%%%%%%%%%%%%%%%%%%%%%%%%%%%%%%%%%%%%%%%%%%
\end{adjustwidth}
\end{document}
